% Generated draft: requires mathematical and source review.
\documentclass[11pt]{article}
\usepackage[margin=25mm]{geometry}
\usepackage{fontspec}
\setmainfont{Latin Modern Roman}
\usepackage{amsmath,amssymb,mathtools,unicode-math}
\setmathfont{Latin Modern Math}
\usepackage{enumitem,needspace,xurl,hyperref}
\setlength{\parindent}{0pt}
\setlength{\parskip}{5pt}
\setlength{\emergencystretch}{3em}
\newcommand{\N}{\mathbb{N}}
\newcommand{\Z}{\mathbb{Z}}
\newcommand{\Q}{\mathbb{Q}}
\newcommand{\R}{\mathbb{R}}
\newcommand{\C}{\mathbb{C}}
\newcommand{\F}{\mathbb{F}}
\newcommand{\A}{\mathbb{A}}
\newcommand{\PP}{\mathbb{P}}
\newcommand{\E}{\mathbb{E}}
\newcommand{\eps}{\varepsilon}
\newcommand{\dd}{\,\mathrm{d}}
\begin{document}
\section*{UnsolvedMath ID: 20003101}
\textbf{Source problem code:} catalogue entry
\subsection{Definitions and mathematical statement}
Let $d_{\mathrm{Th}}$ be the asymmetric Thurston metric on finite-type Teichm\"uller space. A positively unit-speed geodesic ray $X_t$, $t\geq0$, satisfies $d_{\mathrm{Th}}(X_s,X_t)=t-s$ for $s\leq t$. The canonical question asks when two Thurston geodesic rays stay bounded distance apart, but does not specify direction, clocks or comparison of images. These choices matter for an asymmetric metric.

For two rays $X_t,Y_t$, synchronous two-sided boundedness means
\[
 \sup_{t\geq0}\max\{d_{\mathrm{Th}}(X_t,Y_t),
                         d_{\mathrm{Th}}(Y_t,X_t)\}<\infty.
\]
Forward-only or reverse-only bounds omit one term. Reparametrized two-sided fellow traveling instead compares $X_t$ with $Y_{\rho(t)}$ for a specified admissible increasing clock $\rho$, possibly after deleting initial segments. Image-Hausdorff comparison forgets the clocks and requires a choice of symmetrization, for example $d_{\mathrm{sym}}(x,y)=\max\{d_{\mathrm{Th}}(x,y),d_{\mathrm{Th}}(y,x)\}$. The source gives none of these conventions.

The supplied aid treats positively unit-speed cylindrical stretch rays whose directing supports satisfy its chain-recurrence hypotheses. Cylindrical means that the time-zero horocyclic measured foliation is represented by a weighted multicurve. Let those full marked measured multicurves be $\mu_X,\mu_Y$, including their weights. In that stated setting the reported criterion is
\[
 \text{synchronous two-sided boundedness}
 \quad\Longleftrightarrow\quad\mu_X=\mu_Y.
\]
Equality in projective measured-foliation space is insufficient. If $\mu_Y=c\mu_X$ with $c>0$ and $c\ne1$, the aid reports that one of the directed distances has a lower bound of the form $a e^t-b$ for sufficiently large $t$, with $a>0$. The direction and constants depend on the scale convention and are not supplied. The theorem's precise cylinder, completion and chain-recurrence hypotheses and literature reference are missing, so this is not extended to arbitrary rays.

A general clock-rigidity observation is also reported. If both directed distances between $X_t$ and $Y_{\rho(t)}$ are at most $C$ for all relevant $t$, the asymmetric triangle inequalities for unit-speed rays give
\[
 -C-d_{\mathrm{Th}}(Y_0,X_0)
 \leq t-\rho(t)
 \leq C+d_{\mathrm{Th}}(X_0,Y_0).
\]
Thus two-sided bounded fellow traveling forces bounded clock drift. The converse from bounded clock drift to synchronous two-sided boundedness is not asserted: short forward time differences need not bound the reverse direction uniformly along a ray.

The aid uses a cylindrical ray and a shifted tail to distinguish the formulations. A fixed positive time shift is bounded in the forward direction by the shift itself, yet the reverse distance can diverge in the reported cylindrical example; after an appropriate tail reparametrization the paths coincide. Their symmetrized image-Hausdorff distance can nevertheless be finite because only an initial compact segment is omitted. The detailed example and scaling formula are not provided. These distinctions and the special fixed-clock criterion are partial results, not a classification of all Thurston geodesic rays. The trailing dots in the source are layout rather than an additional mathematical condition.
\subsection{Short English statement}
The reported criterion uses identical clocks and both directions of Thurston distance for a special cylindrical family. It requires equality of the weighted multicurves, not just their projective classes. Reparametrized or image-based comparison can behave differently, even for a ray and its tail.
\subsection{Sources}
\begin{itemize}
\item Local UnsolvedMath registry, permanent ID 20003101, source code catalogue entry; \texttt{aim-topology-notes.json}, notes[188].
\item AIM, \emph{Lipschitz metric on Teichmueller space}, Problem 2.11: \url{https://aimath.org/pastworkshops/lipschitzteichproblems.pdf}. The supplied verification reports checking page 2; this draft does not independently repeat it. No proposer or original bibliography is supplied.
\item Cylindrical stretch-ray boundedness and scale-divergence results are invoked in the supplied aid without authors or complete references.
\item Machine-generated research attempt 1: partial result, partially solved, rigor labelled ``proved,'' low novelty confidence. Not peer reviewed; independently verify before citation. Registry updated 14 August 2026.
\end{itemize}
\end{document}
